\documentclass[11pt]{article}

\usepackage[a4paper,margin=28mm]{geometry}
\usepackage{amsmath,amssymb}
\usepackage{booktabs}
\usepackage[hidelinks]{hyperref}

\title{Preconditioning for Distributed Douglas--Rachford Bundle Adjustment}
\author{}
\date{}

\begin{document}
\maketitle

\section{Purpose}

This note summarizes the comparison of camera-coordinate preconditioners for
the distributed Douglas--Rachford splitting (DRS) bundle-adjustment solver.
The practical conclusion is that the implementation should support at least
standard Jacobi scaling and symmetric Ruiz equilibration as selectable modes.
Standard Jacobi remains the conservative default, while symmetric Ruiz is a
competitive alternative that is marginally better in aggregate.

\section{Coordinate scaling}

Let $H$ denote the symmetric camera curvature matrix used to define the DRS
metric.  A positive diagonal matrix $D$ defines scaled coordinates through
\begin{equation}
  x = D^{-1}z,
  \qquad
  \widehat H = D^{-1} H D^{-1}.
\end{equation}
All tested diagonal methods use the same coordinate-transport mechanism.  A
final geometric-mean normalization removes the arbitrary common scale of $D$.

\subsection{Standard Jacobi}

Standard Jacobi scaling uses
\begin{equation}
  D_{jj} = \sqrt{\max(H_{jj},\tau)},
\end{equation}
where $\tau$ is a small relative numerical floor.  This directly equilibrates
the diagonal of $\widehat H$.

\subsection{Fractional Jacobi}

The fractional variants use
\begin{equation}
  D_{jj} = \max(H_{jj},\tau)^q,
  \qquad q\in\{0.25,0.375,0.625\}.
\end{equation}
They test whether weaker or stronger diagonal scaling improves the outer DRS
trajectory.

\subsection{Damped Jacobi}

The tested damped variant uses
\begin{equation}
  D_{jj} = \sqrt{H_{jj} + \lambda},
  \qquad
  \lambda = 10^{-3}\operatorname{median}\{H_{kk}:H_{kk}>0\}.
\end{equation}
The damping protects small diagonal entries, but this particular damping level
did not improve the final objective.

\subsection{Symmetric Ruiz equilibration}

Symmetric Ruiz equilibration iteratively balances the row and column infinity
norms while preserving symmetry.  Starting with $D=I$, each sweep forms
\begin{equation}
  \widehat H = D^{-1}HD^{-1},
  \qquad
  r_i = \lVert \widehat H_{i,:}\rVert_\infty,
\end{equation}
and updates
\begin{equation}
  D \leftarrow D\,\operatorname{diag}(\sqrt{r_i}).
\end{equation}
The experiment used at most ten sweeps and stopped when
$\max_i|\log(\sqrt{r_i})|<10^{-3}$.  Unlike Jacobi, Ruiz uses off-diagonal
coupling information while retaining a diagonal coordinate transform.

\subsection{Pock-inspired variants}

Two additional diagonal variants were tested.  The first uses the square root
of the absolute column sum, and the second uses the square root of the sum of
squared absolute entries.  These are denoted Pock $\alpha=1$ and Pock
$\alpha=0$, respectively, in the result summary.

\section{Experimental protocol}

Each complete cohort contains the same 29 BAL problems, uses 30 clusters, and
runs 90 outer iterations.  Lower objective values are better.  Comparisons use
matched problems and the best objective reached by iteration 90.  The reported
aggregate delta is the geometric mean of the per-problem cost ratios relative
to standard Jacobi.

The explicitly labeled cohorts are stored in
\texttt{serverTest/results\_server.json}.  The standard-Jacobi baseline is the
complete 29-problem unlabeled batch immediately preceding the labeled
\texttt{pock\_gmean} cohort.  Its problem order exactly matches the subsequent
cohorts.  Timing data are available for the five most recent Jacobi and Ruiz
cohorts through the corresponding \texttt{status.tsv} files.

\subsection{Scene normalization by points}

Scene normalization and landmark preconditioning are distinct operations.  At
input, \texttt{normalize\_by\_points} subtracts the coordinate-wise median of
the landmarks and scales the 95th-percentile landmark radius to 100.  Camera
centers are transformed consistently.  This scene normalization is enabled in
the current implementation and was enabled for the preconditioner comparison
above.

An earlier result section explicitly marked
\texttt{normalize\_by\_points off} can be compared with the subsequent
standard-Jacobi section on 28 matched problems.  Enabling normalization lowers
the geometric-mean objective by $0.038\%$ at iteration 30, $0.186\%$ at
iteration 60, and $0.126\%$ at the final checkpoint.  The final comparison has
14 wins for normalization on, 13 wins for normalization off, and one tie.

The large \texttt{problem-3068} instance shows the clearest practical reason to
keep normalization enabled.  Its costs are
\begin{center}
\begin{tabular}{lrrr}
  \toprule
  Setting & Iteration 30 & Iteration 60 & Final \\
  \midrule
  Point normalization off & 3643889 & 3524489 & 3399644 \\
  Point normalization on  & 3646996 & 3384186 & 3299301 \\
  Relative change         & $+0.085\%$ & $-3.981\%$ & $-2.952\%$ \\
  \bottomrule
\end{tabular}
\end{center}
Normalization is initially neutral on this problem but substantially improves
the later trajectory.  It should therefore remain on by default.

Dynamic landmark metric scaling is currently off: the protocol sends unit
\texttt{vnorm} values, so only camera coordinates receive the tested Jacobi or
Ruiz preconditioner.  This setting should not be changed based on the scene
normalization experiment; nontrivial landmark \texttt{vnorm} would require a
separate controlled comparison.

\section{Results}

\begin{table}[ht]
  \centering
  \caption{Final objective relative to standard Jacobi.  A negative delta is
  an improvement.  Runtime is the sum over all 29 problems when recorded.}
  \begin{tabular}{lrrr}
    \toprule
    Method & Geometric cost delta & Cases beating Jacobi & Runtime (s) \\
    \midrule
    Symmetric Ruiz       & $-0.005\%$ & $6/29$ & 8192 \\
    Standard Jacobi      & baseline   & ---    & ---  \\
    Jacobi $q=0.625$     & $+0.563\%$ & $3/29$ & 8276 \\
    Pock $\alpha=0$      & $+0.654\%$ & $2/29$ & ---  \\
    Jacobi $q=0.375$     & $+0.891\%$ & $3/29$ & 8137 \\
    Damped Jacobi        & $+1.167\%$ & $1/29$ & 7948 \\
    Pock $\alpha=1$      & $+1.180\%$ & $3/29$ & ---  \\
    Jacobi $q=0.25$      & $+2.019\%$ & $0/29$ & 7917 \\
    \bottomrule
  \end{tabular}
\end{table}

Symmetric Ruiz and standard Jacobi are nearly tied.  In their direct
comparison, Ruiz wins six problems, Jacobi wins four, and 19 final costs are
identical at the recorded precision.  Ruiz has a geometric-mean improvement of
only $0.005\%$ and a median change of zero.  Its largest improvement is
$0.111\%$ on \texttt{problem-3068}, while its largest regression is $0.043\%$
on \texttt{problem-931}.

At intermediate checkpoints, standard Jacobi has the better mean rank, while
Ruiz records more per-problem wins.  By the final checkpoint, their aggregate
costs have effectively converged.  Ruiz therefore does not replace Jacobi
unconditionally, but it is the only tested alternative that matches or
slightly improves the standard baseline across the complete suite.

Damped Jacobi is about as fast as the other diagonal variants but is worse by
$1.167\%$ in geometric mean and beats standard Jacobi on only one problem.
The fractional and Pock-inspired variants also underperform standard Jacobi in
aggregate.  Of these alternatives, $q=0.625$ is the strongest, but it remains
$0.563\%$ worse geometrically.

\section{Block-Jacobi experiment}

A full $9\times9$ camera-block inverse-square-root transform was implemented as
an optional protocol path.  The original implementation used a relative
eigenvalue floor of $10^{-12}$.  Although the tested camera blocks were finite
and positive definite, their effective condition number reached
$3.65\times10^8$; inversion amplified the weak directions enough that the
first local solve produced non-finite objective values.

The block inverse square root now uses a relative eigenvalue floor of
$10^{-6}$, limiting each regularized block condition number to $10^6$ and the
inverse-square-root condition number to $10^3$.  Coordinate transport also
uses a linear solve instead of explicitly forming the inverse transform.  The
problem-49 smoke test now completes ten outer iterations with finite values.
Its final checkpoint cost was 40660, compared with 37442 for standard Jacobi,
so numerical stability does not yet establish a quality improvement.  A full
block-Jacobi cohort has not been run and no block result is included in the
comparison table.

\section{Implementation conclusion}

The solver should expose at least the following two runtime-selectable camera
preconditioners:
\begin{center}
\begin{tabular}{ll}
  \toprule
  Setting & Meaning \\
  \midrule
  \texttt{jacobi\_gmean} & Standard diagonal Jacobi scaling \\
  \texttt{symmetric\_ruiz\_gmean} & Symmetric Ruiz equilibration \\
  \bottomrule
\end{tabular}
\end{center}
They are currently selected through the environment variable
\texttt{BUNDLE\_PALM\_DRS\_SCALING}.  Standard Jacobi should remain the default
because it is simpler, cheaper to construct, and statistically indistinguishable
from Ruiz on this suite.  Ruiz should remain available for problem-dependent
selection and further experiments because it achieved the best aggregate final
objective and uses coupling information ignored by Jacobi.

\end{document}
